\documentclass[11pt]{memoir}
\usepackage{raeez-math-template}

\title{The Determinant of an Operator}
\author{Raeez Lorgat}
\date{}

\newcommand{\Vect}{\mathrm{Vect}}
\newcommand{\Perf}{\mathrm{Perf}}
\newcommand{\Pic}{\mathrm{Pic}}
\newcommand{\End}{\operatorname{End}}
\newcommand{\Hom}{\operatorname{Hom}}
\newcommand{\Spec}{\operatorname{Spec}}
\newcommand{\Tr}{\operatorname{Tr}}
\newcommand{\Frob}{\operatorname{Frob}}
\newcommand{\rank}{\operatorname{rank}}
\newcommand{\coker}{\operatorname{coker}}
\newcommand{\im}{\operatorname{im}}
\newcommand{\id}{\operatorname{id}}
\newcommand{\cO}{\mathcal{O}}
\newcommand{\cL}{\mathcal{L}}
\newcommand{\cE}{\mathcal{E}}
\newcommand{\cF}{\mathcal{F}}
\newcommand{\cD}{\mathcal{D}}
\newcommand{\cH}{\mathcal{H}}
\newcommand{\bA}{\mathbb{A}}
\newcommand{\bC}{\mathbb{C}}
\newcommand{\bF}{\mathbb{F}}
\newcommand{\bQ}{\mathbb{Q}}
\newcommand{\bR}{\mathbb{R}}
\newcommand{\bZ}{\mathbb{Z}}

\begin{document}

\maketitle
\tableofcontents

\chapter*{North Star}
\addcontentsline{toc}{chapter}{North Star}

The determinant of an operator is the scalar by which the operator acts
on the top exterior line.

Let \(V\) be an \(n\)-dimensional vector space over a field \(k\). The
top exterior power
\[
  \det V := \bigwedge^{n}V
\]
is one-dimensional. A linear operator \(T:V\to V\) induces a linear
operator
\[
  \bigwedge^{n}T:\bigwedge^{n}V\longrightarrow \bigwedge^{n}V.
\]
Every endomorphism of a one-dimensional vector space is multiplication
by a unique scalar. That scalar is \(\det(T)\):
\[
  \bigwedge^{n}T=\det(T)\cdot \id_{\det V}.
\]

This is the definition. The formula with permutations, the product of
eigenvalues, the Jacobian, the Frobenius determinant in a zeta
function, the determinant line on a moduli stack, the Weyl denominator,
and the central extension of a loop group are realizations of this
definition.

\chapter{The Finite-Dimensional Construction}
\label{ch:finite}

\section{Alternating Volume Before Determinants}
\label{sec:alternating-volume}

The determinant is forced by one elementary fact: volume is alternating.
If two edge vectors of a parallelepiped are equal, the parallelepiped
has zero volume. If two edge vectors are exchanged, the orientation
changes sign. Algebra encodes this by the exterior power.

\begin{definition}[alternating multilinear map]
\label{def:alternating-map}
Let \(V\) and \(L\) be vector spaces over a field \(k\). A map
\[
  A:V^r\to L
\]
is alternating \(r\)-linear if it is linear in each variable and
\[
  A(v_1,\dots,v_r)=0
\]
whenever \(v_i=v_j\) for some \(i\neq j\).
\end{definition}

From this one derives the sign rule. If \(A\) is alternating and
\(i\neq j\), then
\[
  A(\dots,v_i,\dots,v_j,\dots)
  =
  -A(\dots,v_j,\dots,v_i,\dots).
\]
Indeed, in the two displayed slots,
\[
  0=A(\dots,v_i+v_j,\dots,v_i+v_j,\dots),
\]
and multilinearity expands this as the sum of the two swapped terms,
because the repeated terms vanish.

\begin{definition}[exterior power]
\label{def:exterior-power}
The \(r\)-th exterior power \(\bigwedge^rV\) is the vector space
generated by symbols
\[
  v_1\wedge\cdots\wedge v_r
\]
subject to the relations making the map
\[
  V^r\to \bigwedge^rV,\qquad
  (v_1,\dots,v_r)\mapsto v_1\wedge\cdots\wedge v_r
\]
alternating and multilinear.
\end{definition}

Equivalently, \(\bigwedge^rV\) is the quotient of \(V^{\otimes r}\)
by the subspace generated by all tensors which express failure of
alternation. The quotient formulation proves existence. The symbolic
formulation proves usability.

\begin{proposition}[universal property of exterior powers]
\label{prop:exterior-universal}
For every vector space \(L\), composition with the wedge map gives a
natural bijection
\[
  \Hom\left(\bigwedge^rV,L\right)
  \simeq
  \{\text{alternating \(r\)-linear maps }V^r\to L\}.
\]
\end{proposition}

\begin{proof}
A linear map \(\phi:\bigwedge^rV\to L\) gives an alternating
\(r\)-linear map
\[
  (v_1,\dots,v_r)\mapsto \phi(v_1\wedge\cdots\wedge v_r).
\]
Conversely, if \(A:V^r\to L\) is alternating and multilinear, then
the defining relations of \(\bigwedge^rV\) are exactly the relations
that \(A\) already satisfies. Hence \(A\) descends uniquely to a
linear map \(\phi_A:\bigwedge^rV\to L\) with
\[
  \phi_A(v_1\wedge\cdots\wedge v_r)=A(v_1,\dots,v_r).
\]
The two constructions are inverse.
\end{proof}

\begin{proposition}[basis of an exterior power]
\label{prop:wedge-basis}
Let \(V\) have basis \(e_1,\dots,e_n\). Then \(\bigwedge^rV\) has
basis
\[
  e_{i_1}\wedge\cdots\wedge e_{i_r},
  \qquad
  1\leq i_1<\cdots<i_r\leq n.
\]
In particular,
\[
  \dim\bigwedge^rV=\binom nr.
\]
\end{proposition}

\begin{proof}
First, the displayed vectors span. Every wedge
\[
  v_1\wedge\cdots\wedge v_r
\]
expands multilinearly in the basis \(e_i\). Any term with repeated
indices vanishes. Any term with distinct indices can be reordered
into increasing order, at the cost of the sign of the permutation.

Second, the displayed vectors are independent. For each increasing
multi-index \(I=(i_1<\cdots<i_r)\), define an alternating \(r\)-linear
form \(\varepsilon_I:V^r\to k\) by
\[
  \varepsilon_I(e_{j_1},\dots,e_{j_r})
  =
  \begin{cases}
    \operatorname{sgn}(\sigma),&
    (j_1,\dots,j_r)=(i_{\sigma(1)},\dots,i_{\sigma(r)})
    \text{ for some }\sigma\in S_r,\\
    0,&\text{otherwise.}
  \end{cases}
\]
By \Cref{prop:exterior-universal}, \(\varepsilon_I\) descends to a
linear functional on \(\bigwedge^rV\). It takes value \(1\) on
\(e_{i_1}\wedge\cdots\wedge e_{i_r}\) and value \(0\) on every other
increasing basis candidate. These functionals separate the displayed
vectors. Hence the spanning family is independent.
\end{proof}

\begin{corollary}[the top exterior power is a line]
\label{cor:top-wedge-line}
If \(\dim V=n\), then
\[
  \bigwedge^nV
\]
is one-dimensional. If \(e_1,\dots,e_n\) is a basis of \(V\), then
\[
  e_1\wedge\cdots\wedge e_n
\]
is a basis of \(\bigwedge^nV\).
\end{corollary}

\begin{proof}
Apply \Cref{prop:wedge-basis} with \(r=n\). There is exactly one
increasing \(n\)-tuple:
\[
  (1,2,\dots,n).
\]
\end{proof}

\section{Linear Maps Act on Exterior Powers}
\label{sec:linear-maps-exterior}

Let \(T:V\to W\) be linear. Define
\[
  \bigwedge^rT:\bigwedge^rV\to\bigwedge^rW
\]
by
\[
  (\bigwedge^rT)(v_1\wedge\cdots\wedge v_r)
  :=
  T(v_1)\wedge\cdots\wedge T(v_r).
\]
This is well-defined because the right side is alternating and
multilinear in \(v_1,\dots,v_r\). The universal property of the
exterior power is exactly the statement that this prescription
descends to a linear map.

\begin{proposition}[functoriality]
\label{prop:wedge-functoriality}
If \(V\xrightarrow{T}W\xrightarrow{S}U\) are linear maps, then
\[
  \bigwedge^r(S\circ T)=(\bigwedge^rS)\circ(\bigwedge^rT).
\]
Also
\[
  \bigwedge^r(\id_V)=\id_{\bigwedge^rV}.
\]
\end{proposition}

\begin{proof}
Both identities are checked on wedge generators. For example,
\[
  \bigwedge^r(S\circ T)(v_1\wedge\cdots\wedge v_r)
  =
  S(Tv_1)\wedge\cdots\wedge S(Tv_r),
\]
which is exactly
\[
  (\bigwedge^rS)(Tv_1\wedge\cdots\wedge Tv_r)
  =
  ((\bigwedge^rS)\circ(\bigwedge^rT))(v_1\wedge\cdots\wedge v_r).
\]
\end{proof}

\section{Definition of the Determinant}
\label{sec:det-definition}

Now take \(W=V\), \(r=n=\dim V\), and \(T:V\to V\). Functoriality
gives
\[
  \bigwedge^nT:\bigwedge^nV\to\bigwedge^nV.
\]
By \Cref{cor:top-wedge-line}, the source and target form a
one-dimensional line. Hence the operator is multiplication by one
scalar.

\begin{definition}[determinant of an operator]
\label{def:operator-det}
Let \(V\) be finite-dimensional and \(T:V\to V\) linear. The
determinant \(\det(T)\in k\) is the unique scalar satisfying
\[
  \bigwedge^{\dim V}T
  =
  \det(T)\cdot \id_{\bigwedge^{\dim V}V}.
\]
\end{definition}

This definition is the north star. The rest of finite-dimensional
determinant theory is obtained by evaluating this scalar in useful
coordinates.

\section{One Dimension}
\label{sec:one-dimension}

Let \(V\) be one-dimensional. Choose a non-zero vector \(e\). Every
linear operator \(T:V\to V\) has the form
\[
  T(e)=a e
\]
for a unique scalar \(a\in k\). In one dimension the determinant is
this scalar:
\[
  \det(T)=a.
\]

This case already contains the whole idea. A one-dimensional vector
space is a line. An operator on a line is multiplication by a scalar.
The determinant asks every operator, even an operator on a large
space, to act on a line canonically attached to that space.

\section{Two Dimensions by Area}
\label{sec:two-dimensions-area}

Let \(V\) have basis \(e_{1},e_{2}\). Write
\[
  T(e_{1})=a e_{1}+c e_{2},\qquad
  T(e_{2})=b e_{1}+d e_{2}.
\]
The matrix of \(T\) in the basis \((e_{1},e_{2})\) is
\[
  [T]=
  \begin{pmatrix}
    a & b\\
    c & d
  \end{pmatrix}.
\]
The oriented area element is \(e_{1}\wedge e_{2}\). The wedge product
is bilinear and alternating:
\[
  u\wedge v=-v\wedge u,\qquad u\wedge u=0.
\]
Therefore
\begin{align*}
  T(e_{1})\wedge T(e_{2})
  &=(a e_{1}+c e_{2})\wedge(b e_{1}+d e_{2})\\
  &=ab\,e_{1}\wedge e_{1}
    +ad\,e_{1}\wedge e_{2}
    +cb\,e_{2}\wedge e_{1}
    +cd\,e_{2}\wedge e_{2}\\
  &=ad\,e_{1}\wedge e_{2}-bc\,e_{1}\wedge e_{2}\\
  &=(ad-bc)e_{1}\wedge e_{2}.
\end{align*}
Hence
\[
  \det(T)=ad-bc.
\]

Nothing was memorized. The formula is forced by the action on the
oriented area line \(\bigwedge^{2}V\).

\section{Three Dimensions by Volume}
\label{sec:three-dimensions-volume}

Let \(V\) have basis \(e_{1},e_{2},e_{3}\). The oriented volume line is
\[
  \det V=\bigwedge^{3}V=k\cdot(e_{1}\wedge e_{2}\wedge e_{3}).
\]
If
\[
  T(e_j)=\sum_{i=1}^{3}a_{ij}e_i,
\]
then
\[
  T(e_1)\wedge T(e_2)\wedge T(e_3)
  =
  \left(\sum_{\sigma\in S_3}\operatorname{sgn}(\sigma)
  a_{\sigma(1),1}a_{\sigma(2),2}a_{\sigma(3),3}\right)
  e_1\wedge e_2\wedge e_3.
\]
Terms with repeated basis vectors vanish because the wedge product is
alternating. Terms using each basis vector once survive. The sign is
the sign of the permutation needed to reorder the surviving wedge into
\(e_1\wedge e_2\wedge e_3\).

Thus the determinant is the signed sum over all ways of selecting one
entry from each column and one entry from each row.

\section{The General Formula}
\label{sec:general-formula}

Let \(V\) have basis \(e_1,\dots,e_n\). Let \(T(e_j)=\sum_i a_{ij}e_i\).
Then
\[
  T(e_1)\wedge\cdots\wedge T(e_n)
  =
  \det(a_{ij})\,e_1\wedge\cdots\wedge e_n,
\]
where
\[
  \det(a_{ij})
  =
  \sum_{\sigma\in S_n}\operatorname{sgn}(\sigma)
  a_{\sigma(1),1}\cdots a_{\sigma(n),n}.
\]

\begin{proof}
Expand the wedge product by multilinearity:
\[
  \left(\sum_{i_1}a_{i_1,1}e_{i_1}\right)\wedge\cdots\wedge
  \left(\sum_{i_n}a_{i_n,n}e_{i_n}\right).
\]
If two indices \(i_r,i_s\) coincide, the term contains the same basis
vector twice and vanishes. The only surviving terms are indexed by
permutations \(\sigma\in S_n\), with \(i_j=\sigma(j)\). For such a
term,
\[
  e_{\sigma(1)}\wedge\cdots\wedge e_{\sigma(n)}
  =
  \operatorname{sgn}(\sigma)\,e_1\wedge\cdots\wedge e_n.
\]
Substitution gives the displayed formula.
\end{proof}

\section{Basis Independence}
\label{sec:basis-independence}

The determinant is not the permutation formula. The formula depends on
a basis. The determinant does not.

\begin{proposition}[basis-free determinant]
\label{prop:basis-free-det}
Let \(V\) be finite-dimensional. There is a unique scalar \(\det(T)\)
such that
\[
  \bigwedge^{\dim V}T=\det(T)\cdot\id_{\det V}.
\]
This scalar is independent of every choice of basis.
\end{proposition}

\begin{proof}
The line \(\det V=\bigwedge^{n}V\) is one-dimensional. Hence every
endomorphism of \(\det V\) is multiplication by a unique scalar.
The induced map \(\bigwedge^n T\) is such an endomorphism. The scalar
is defined without choosing a basis. If a basis is chosen, the scalar
is computed by the permutation formula of
\Cref{sec:general-formula}; changing basis changes both the chosen
generator of \(\det V\) and its image by the same non-zero scalar, so
the multiplier is unchanged.
\end{proof}

\section{What the Determinant Measures}
\label{sec:what-det-measures}

\begin{proposition}[detecting invertibility]
\label{prop:det-invertibility}
For \(T:V\to V\) with \(V\) finite-dimensional,
\[
  \det(T)\neq 0
  \quad\Longleftrightarrow\quad
  T\text{ is invertible}.
\]
\end{proposition}

\begin{proof}
If \(T\) is not invertible, its image has dimension \(<n\). Then
\(T(e_1),\dots,T(e_n)\) are linearly dependent for every basis
\((e_i)\). The wedge of linearly dependent vectors is zero, so
\(\bigwedge^nT=0\), hence \(\det(T)=0\).

If \(T\) is invertible, \(T(e_1),\dots,T(e_n)\) is a basis. Its wedge
is a non-zero generator of \(\det V\). Hence
\(\bigwedge^nT\) is non-zero on a one-dimensional space, so it is
multiplication by a non-zero scalar.
\end{proof}

\begin{proposition}[multiplicativity]
\label{prop:multiplicativity}
For \(S,T:V\to V\),
\[
  \det(ST)=\det(S)\det(T).
\]
\end{proposition}

\begin{proof}
Exterior powers are functorial:
\[
  \bigwedge^n(ST)=(\bigwedge^nS)(\bigwedge^nT).
\]
On the determinant line, these maps are multiplication by
\(\det(S)\) and \(\det(T)\). Their composite is multiplication by
\(\det(S)\det(T)\).
\end{proof}

\begin{proposition}[eigenvalue product]
\label{prop:eigenvalue-product}
If \(T\) has eigenvalues \(\lambda_1,\dots,\lambda_n\) counted with
algebraic multiplicity over an algebraic closure of \(k\), then
\[
  \det(T)=\lambda_1\cdots\lambda_n.
\]
\end{proposition}

\begin{proof}
Over an algebraic closure, triangularize \(T\). In an upper triangular
matrix, the determinant formula has only one non-zero surviving
diagonal contribution:
\[
  \det(T)=a_{11}\cdots a_{nn}.
\]
The diagonal entries of a triangular form are the eigenvalues counted
with algebraic multiplicity. The determinant is unchanged by extension
of scalars.
\end{proof}

\section{Worked Linear Examples}
\label{sec:worked-linear-examples}

\subsection{Diagonal Scaling}

Let \(V=k^3\) with basis \(e_1,e_2,e_3\), and let
\[
  T(e_1)=a e_1,\qquad T(e_2)=b e_2,\qquad T(e_3)=c e_3.
\]
Then
\[
  T(e_1)\wedge T(e_2)\wedge T(e_3)
  =
  (a e_1)\wedge(b e_2)\wedge(c e_3)
  =
  abc\,e_1\wedge e_2\wedge e_3.
\]
Thus
\[
  \det(T)=abc.
\]
The determinant is the product of the three independent stretch
factors.

\subsection{Shear}

Let \(V=k^2\), and let
\[
  T(e_1)=e_1,\qquad T(e_2)=e_2+a e_1.
\]
Then
\[
  T(e_1)\wedge T(e_2)
  =
  e_1\wedge(e_2+a e_1)
  =
  e_1\wedge e_2+a\,e_1\wedge e_1
  =
  e_1\wedge e_2.
\]
Therefore
\[
  \det(T)=1.
\]
A shear changes shape but not oriented area.

\subsection{Projection}

Let \(V=k^2\), and let
\[
  P(e_1)=e_1,\qquad P(e_2)=0.
\]
Then
\[
  P(e_1)\wedge P(e_2)=e_1\wedge0=0.
\]
Hence
\[
  \det(P)=0.
\]
Projection collapses area to a line. The determinant vanishes because
the top exterior image vanishes.

\subsection{\texorpdfstring{Rotation over \(\bR\)}{Rotation over R}}

Let \(V=\bR^2\), and let
\[
  R_\theta=
  \begin{pmatrix}
    \cos\theta & -\sin\theta\\
    \sin\theta & \cos\theta
  \end{pmatrix}.
\]
Then
\[
  R_\theta(e_1)=\cos\theta\,e_1+\sin\theta\,e_2,
  \qquad
  R_\theta(e_2)=-\sin\theta\,e_1+\cos\theta\,e_2.
\]
Therefore
\begin{align*}
  R_\theta(e_1)\wedge R_\theta(e_2)
  &=
  (\cos\theta\,e_1+\sin\theta\,e_2)
  \wedge
  (-\sin\theta\,e_1+\cos\theta\,e_2)\\
  &=
  \cos^2\theta\,e_1\wedge e_2
  +
  \sin^2\theta\,e_1\wedge e_2\\
  &=
  e_1\wedge e_2.
\end{align*}
Thus
\[
  \det(R_\theta)=1.
\]
Rotation preserves oriented area.

\section{Exact Sequences and Block Operators}
\label{sec:exact-sequences}

The determinant line behaves well in exact sequences. This is the
first point at which the finite-dimensional definition becomes a
machine for geometry.

\begin{proposition}[determinant line of a short exact sequence]
\label{prop:short-exact-det-line}
Let
\[
  0\longrightarrow U\longrightarrow V\longrightarrow W\longrightarrow 0
\]
be a short exact sequence of finite-dimensional vector spaces. There
is a canonical isomorphism
\[
  \det V\simeq \det U\otimes \det W.
\]
\end{proposition}

\begin{proof}
Choose a basis \(u_1,\dots,u_r\) of \(U\). Choose vectors
\(\widetilde w_1,\dots,\widetilde w_s\in V\) lifting a basis
\(w_1,\dots,w_s\) of \(W\). Then
\[
  u_1\wedge\cdots\wedge u_r\wedge
  \widetilde w_1\wedge\cdots\wedge \widetilde w_s
\]
is a basis of \(\det V\). Send it to
\[
  (u_1\wedge\cdots\wedge u_r)\otimes
  (w_1\wedge\cdots\wedge w_s).
\]
Changing the lifts \(\widetilde w_i\) adds \(U\)-vectors to them; the
extra terms contain repeated \(U\)-directions in the top wedge and
vanish. Changing bases multiplies both sides by the same determinant.
Hence the map is canonical.
\end{proof}

\begin{corollary}[block triangular determinant]
\label{cor:block-triangular}
If \(T:V\to V\) preserves a subspace \(U\subset V\), and if
\(\overline T\) is the induced operator on \(V/U\), then
\[
  \det(T)=\det(T|_U)\det(\overline T).
\]
\end{corollary}

\begin{proof}
Apply \(\det\) to the exact sequence
\[
  0\to U\to V\to V/U\to 0
\]
and use \Cref{prop:short-exact-det-line}. On the line
\(\det U\otimes\det(V/U)\), the induced action is multiplication by
\(\det(T|_U)\det(\overline T)\).
\end{proof}

\chapter{From Operators to Determinant Lines}
\label{ch:det-lines}

\section{The Determinant of a Vector Bundle}
\label{sec:vector-bundle-det}

Let \(X\) be a scheme or a smooth manifold. A rank \(n\) vector bundle
\(E\) assigns an \(n\)-dimensional vector space \(E_x\) to every point
\(x\in X\), varying algebraically or smoothly. Its determinant is the
line bundle
\[
  \det E:=\bigwedge^n E.
\]
Fibrewise,
\[
  (\det E)_x=\det(E_x).
\]

Thus the finite-dimensional determinant line globalizes without
change. A vector bundle is a family of vector spaces. Its determinant
line is the family of their top exterior lines.

\section{Determinant of a Bundle Endomorphism}
\label{sec:bundle-endomorphism}

Let \(T:E\to E\) be an endomorphism of a rank \(n\) vector bundle.
Taking top exterior powers gives
\[
  \bigwedge^nT:\det E\to\det E.
\]
An endomorphism of a line bundle is a section of \(\cO_X\). Thus
\[
  \det(T)\in \Gamma(X,\cO_X).
\]
At a point \(x\),
\[
  \det(T)(x)=\det(T_x:E_x\to E_x).
\]
The zero locus of this function is exactly the locus where \(T_x\) is
not invertible.

\section{The Determinant of a Complex}
\label{sec:det-complex}

Geometry rarely gives a single vector space. It gives cohomology in
several degrees. The determinant line of a complex is the alternating
product of the determinant lines of its cohomology.

\begin{definition}[determinant line of a finite complex]
\label{def:det-complex}
Let \(C^\bullet\) be a bounded complex of finite-dimensional vector
spaces. Define
\[
  \det C^\bullet
  :=
  \bigotimes_i \left(\det H^i(C^\bullet)\right)^{(-1)^i},
\]
where \(L^{+1}=L\) and \(L^{-1}=L^\vee\).
\end{definition}

The inverse in odd degree is not decoration. It is the line-level
form of the alternating sign in Euler characteristics:
\[
  \chi(C^\bullet)=\sum_i(-1)^i\dim H^i(C^\bullet).
\]
The determinant line is the multiplicative Euler characteristic.

\begin{example}[two-term complex]
\label{ex:two-term-complex}
Let
\[
  C^\bullet=[C^0\xrightarrow{d}C^1].
\]
Then
\[
  \det C^\bullet
  =
  \det H^0(C^\bullet)\otimes \det H^1(C^\bullet)^\vee.
\]
If \(d\) is an isomorphism, both cohomology groups vanish and
\(\det C^\bullet\simeq k\). Acyclic complexes have canonically
trivial determinant lines after choosing the differential. This is
the origin of torsion invariants.
\end{example}

\section{Why the Signs Are Forced}
\label{sec:signs-forced}

Suppose \(C^\bullet\) is a complex with zero differential:
\[
  C^\bullet=\cdots\to 0\to C^0\to C^1\to 0\to\cdots.
\]
Then \(H^i(C^\bullet)=C^i\), so
\[
  \det C^\bullet=\det C^0\otimes(\det C^1)^\vee.
\]
If one thinks of \(C^0\) as numerator and \(C^1\) as denominator, then
the determinant line remembers the cancellation encoded by the
differential. This is the same alternating structure as
\[
  \det(1-tT|C^\bullet)
  :=
  \prod_i\det(1-tT|C^i)^{(-1)^i}.
\]
The signs are the price of making cohomology invariant under
homotopy.

\chapter{Infinite-Dimensional Operators}
\label{ch:infinite}

\section{The Obstruction}
\label{sec:infinite-obstruction}

The finite-dimensional determinant used the line
\(\bigwedge^nV\). If \(V\) is infinite-dimensional, there is no final
exterior power. There is no largest \(n\). Thus there is no automatic
line \(\det V\), and no automatic scalar \(\det(T)\).

Every infinite-dimensional determinant is therefore extra structure.
One must specify:
\begin{enumerate}
\item a class of operators for which a determinant is defined;
\item a replacement for the finite product of eigenvalues;
\item a topology or regularisation making the replacement converge;
\item functorial identities preserved by the construction.
\end{enumerate}

\section{Fredholm Determinants}
\label{sec:fredholm}

Let \(\cH\) be a separable Hilbert space and let \(K:\cH\to\cH\) be
trace-class. The operator \(1+K\) has a Fredholm determinant
\[
  \det(1+K).
\]
If \(K\) has finite rank with eigenvalues \(\lambda_1,\dots,\lambda_m\)
on its non-zero image, then
\[
  \det(1+K)=\prod_{j=1}^{m}(1+\lambda_j).
\]
For trace-class \(K\), the infinite product over eigenvalues is
convergent after the trace-class hypothesis is imposed:
\[
  \det(1+K)=\prod_j(1+\lambda_j).
\]
Equivalently,
\[
  \det(1+K)
  =
  \sum_{r\geq0}\Tr\left(\bigwedge^rK\right),
\]
where the series converges for trace-class \(K\). This formula is the
finite-dimensional exterior-power construction with all exterior
degrees retained, because there is no top degree.

\section{Zeta-Regularised Determinants}
\label{sec:zeta-det}

Let \(A\) be a positive self-adjoint operator with discrete positive
spectrum
\[
  0<\lambda_1\leq\lambda_2\leq\cdots
\]
and finite multiplicities. Suppose the spectral zeta function
\[
  \zeta_A(s):=\sum_j\lambda_j^{-s}
\]
converges for \(\operatorname{Re}(s)\gg0\) and admits meromorphic
continuation to a neighbourhood of \(s=0\), regular at \(0\). Define
\[
  \det_\zeta(A):=\exp\bigl(-\zeta_A'(0)\bigr).
\]

This definition is forced by the finite case. If the spectrum is
finite, then
\[
  \zeta_A'(s)=\sum_j -(\log\lambda_j)\lambda_j^{-s},
\]
so
\[
  \zeta_A'(0)=-\sum_j\log\lambda_j.
\]
Therefore
\[
  \exp(-\zeta_A'(0))
  =
  \exp\left(\sum_j\log\lambda_j\right)
  =
  \prod_j\lambda_j.
\]
The zeta determinant is the finite product of eigenvalues continued
through the spectral zeta function.

\section{\texorpdfstring{The Determinant of \(s-\Theta\)}{The Determinant of s-Theta}}
\label{sec:s-minus-theta}

In arithmetic applications one often wants a determinant depending on
a complex parameter \(s\):
\[
  \det_\infty(s\id-\Theta).
\]
If \(\Theta\) has discrete spectrum \(\{\rho\}\), this determinant is
morally
\[
  \prod_\rho(s-\rho).
\]
Such a product usually diverges. A regularised determinant is a rule
which makes sense of this product while preserving the expected
logarithmic derivative:
\[
  \frac{d}{ds}\log\det_\infty(s\id-\Theta)
  =
  \Tr_{\mathrm{reg}}\bigl((s\id-\Theta)^{-1}\bigr).
\]
Thus a determinant packages a spectrum into a function. Conversely,
if a function is known by its zeros, a determinant formula asserts
that those zeros are the spectrum of an operator.

This is why determinant formulas in arithmetic are severe. Writing
\[
  \xi(s)=\det_\infty(s\id-\Theta)
\]
does not merely restate the zero set of \(\xi\). It asserts that the
zeros are realised by an operator \(\Theta\) on a cohomological space,
with the determinant regularisation compatible with trace formulas.

\chapter{Algebraic Geometry}
\label{ch:algebraic-geometry}

\section{The Jacobian as a Determinant}
\label{sec:jacobian}

Let \(f:X\to Y\) be a morphism of smooth \(n\)-dimensional varieties
over a field \(k\). The differential is a bundle map
\[
  df:f^*\Omega^1_Y\to\Omega^1_X.
\]
Taking determinants gives
\[
  \det(df):
  f^*\det\Omega^1_Y\longrightarrow \det\Omega^1_X.
\]
Locally, if \(Y\) has coordinates \(y_1,\dots,y_n\) and \(X\) has
coordinates \(x_1,\dots,x_n\), then
\[
  df(dy_j)=\sum_i\frac{\partial f_j}{\partial x_i}dx_i.
\]
Therefore
\[
  df(dy_1)\wedge\cdots\wedge df(dy_n)
  =
  \det\left(\frac{\partial f_j}{\partial x_i}\right)
  dx_1\wedge\cdots\wedge dx_n.
\]
The classical Jacobian determinant is the determinant of the
differential on the canonical line.

\section{Canonical Bundle}
\label{sec:canonical-bundle}

For a smooth \(n\)-dimensional variety \(X\), the canonical bundle is
\[
  \omega_X:=\det\Omega^1_X=\bigwedge^n\Omega^1_X.
\]
This is the line of algebraic volume forms. The determinant
construction is the reason the canonical bundle is a line: top
forms are top exterior powers of cotangent spaces.

If \(f:X\to Y\) is etale of relative dimension \(0\), then \(df\) is
an isomorphism. Hence \(\det(df)\) is non-zero everywhere. Conversely,
for a finite separable map between smooth curves, vanishing of
\(\det(df)\) detects ramification. In dimension \(1\), this is the
ordinary derivative. In higher dimension, it is the determinant of the
Jacobian matrix.

\section{Chern Classes and the Determinant}
\label{sec:chern-det}

For a vector bundle \(E\) of rank \(n\), the first Chern class is the
first Chern class of its determinant line:
\[
  c_1(E)=c_1(\det E).
\]
The determinant extracts the line-bundle shadow of \(E\). Many
intersection-theoretic invariants of \(E\) begin with this shadow.

If \(E=L_1\oplus\cdots\oplus L_n\) splits as line bundles, then
\[
  \det E=L_1\otimes\cdots\otimes L_n,
\]
so
\[
  c_1(E)=c_1(L_1)+\cdots+c_1(L_n).
\]
This is the line-bundle version of the eigenvalue product formula.

\section{Determinant of Cohomology}
\label{sec:det-cohomology-ag}

Let \(X\) be a proper scheme over \(k\), and let \(E\) be a coherent
sheaf whose cohomology groups are finite-dimensional. Define
\[
  \lambda(E)
  :=
  \det R\Gamma(X,E)
  =
  \bigotimes_i \det H^i(X,E)^{(-1)^i}.
\]
This line is the determinant of cohomology.

For a smooth projective curve \(C\),
\[
  \lambda(E)=\det H^0(C,E)\otimes \det H^1(C,E)^\vee.
\]
If \(E\) varies in a family, these lines form a line bundle on the
parameter space. Theta line bundles on moduli spaces are determinant
of cohomology line bundles.

\section{Degeneracy Loci}
\label{sec:degeneracy}

Let \(\phi:E\to F\) be a morphism of vector bundles of equal rank
\(n\) on \(X\). Taking determinants gives
\[
  \det(\phi):\det E\to\det F.
\]
Equivalently,
\[
  \det(\phi)\in\Gamma(X,\det F\otimes(\det E)^\vee).
\]
The vanishing locus of this section is the locus where \(\phi_x\) is
not an isomorphism. Thus a determinant turns a rank condition into a
divisor when the expected codimension is \(1\).

More general degeneracy loci use minors. A minor is the determinant of
a map between exterior powers:
\[
  \bigwedge^r E\to \bigwedge^r F.
\]
Thus even the higher rank stratification is built from determinant
maps.

\section{Derived Algebraic Geometry}
\label{sec:derived-ag-det}

In derived algebraic geometry the basic objects are perfect complexes,
not only vector bundles. The determinant extends to a functor
\[
  \det:\Perf(X)\to \Pic(X),
\]
where \(\Pic(X)\) is the Picard groupoid of line bundles. Locally, if
a perfect complex is represented by vector bundles \(E^i\), then
\[
  \det(E^\bullet)=\bigotimes_i(\det E^i)^{(-1)^i}.
\]
Passing to cohomology gives the same line when the complex has finite
cohomology. This determinant functor is the line-level shadow of the
derived category.

The point is structural. Algebraic geometry uses determinants because
geometry produces families. A determinant is the functor which turns a
family of finite-dimensional linear problems into a line bundle, and
line bundles have divisors, Chern classes, metrics, and moduli.

\chapter{Arithmetic Geometry}
\label{ch:arithmetic-geometry}

\section{The Fundamental Trace-to-Determinant Identity}
\label{sec:trace-det-identity}

Let \(T\) be an endomorphism of a finite-dimensional vector space.
Then
\[
  -\log\det(1-tT)
  =
  \sum_{r\geq1}\frac{\Tr(T^r)}{r}t^r
\]
as a formal power series.

\begin{proof}
After extending scalars, triangularize \(T\), with eigenvalues
\(\lambda_1,\dots,\lambda_n\). Then
\[
  \det(1-tT)=\prod_j(1-t\lambda_j).
\]
Hence
\[
  -\log\det(1-tT)
  =
  -\sum_j\log(1-t\lambda_j)
  =
  \sum_j\sum_{r\geq1}\frac{(t\lambda_j)^r}{r}.
\]
Interchange the finite sum over \(j\) with the formal sum over \(r\):
\[
  \sum_{r\geq1}\frac{t^r}{r}\sum_j\lambda_j^r.
\]
But \(\sum_j\lambda_j^r=\Tr(T^r)\). This proves the identity.
\end{proof}

This identity is the bridge from point counts to zeta functions.

\section{Varieties over Finite Fields}
\label{sec:finite-field-zeta}

Let \(X/\bF_q\) be a variety. Its zeta function is
\[
  Z(X,t):=
  \exp\left(\sum_{r\geq1}\#X(\bF_{q^r})\frac{t^r}{r}\right).
\]
The Grothendieck-Lefschetz trace formula says
\[
  \#X(\bF_{q^r})
  =
  \sum_i(-1)^i
  \Tr\left(\Frob_q^r\mid H^i_{\mathrm{\acute et}}(X_{\overline{\bF}_q},\bQ_\ell)\right).
\]
Insert this into the exponential defining \(Z(X,t)\):
\[
  \log Z(X,t)
  =
  \sum_i(-1)^i
  \sum_{r\geq1}
  \Tr\left(\Frob_q^r\mid H^i\right)\frac{t^r}{r}.
\]
By \Cref{sec:trace-det-identity},
\[
  \sum_{r\geq1}\Tr(\Frob_q^r|H^i)\frac{t^r}{r}
  =
  -\log\det(1-t\Frob_q|H^i).
\]
Therefore
\[
  Z(X,t)
  =
  \prod_i
  \det\left(1-t\Frob_q\mid
  H^i_{\mathrm{\acute et}}(X_{\overline{\bF}_q},\bQ_\ell)\right)^{(-1)^{i+1}}.
\]

This is the arithmetic determinant formula. Counting points over all
finite extensions is equivalent to taking alternating determinants of
Frobenius on cohomology.

\section{Example: A Smooth Projective Curve}
\label{sec:curve-zeta}

Let \(C/\bF_q\) be a smooth projective geometrically connected curve
of genus \(g\). Its etale cohomology has three non-zero degrees:
\[
  H^0\simeq \bQ_\ell,\qquad
  \dim H^1=2g,\qquad
  H^2\simeq \bQ_\ell(-1).
\]
Frobenius acts on \(H^0\) by \(1\), and on \(H^2\) by \(q\). Hence
\[
  Z(C,t)
  =
  \frac{\det(1-t\Frob_q|H^1)}
  {(1-t)(1-qt)}.
\]
Usually one writes
\[
  Z(C,t)=\frac{P_C(t)}{(1-t)(1-qt)},
\qquad
  P_C(t)=\det(1-t\Frob_q|H^1).
\]
The numerator is a determinant on middle cohomology. The denominator
comes from the two Tate lines \(H^0\) and \(H^2\).

\section{Example: An Elliptic Curve}
\label{sec:elliptic-zeta}

Let \(E/\bF_q\) be an elliptic curve. Then \(g=1\), so \(H^1\) has
dimension \(2\). Let
\[
  a_q:=q+1-\#E(\bF_q).
\]
The Frobenius characteristic polynomial on \(H^1\) is
\[
  1-a_qt+qt^2.
\]
Therefore
\[
  Z(E,t)=\frac{1-a_qt+qt^2}{(1-t)(1-qt)}.
\]
The Riemann hypothesis for \(E/\bF_q\) says the two reciprocal roots
of \(1-a_qt+qt^2\) have absolute value \(q^{1/2}\). In determinant
language, it is a purity statement about Frobenius eigenvalues on
\(H^1\).

\section{\texorpdfstring{Local \(L\)-Factors}{Local L-Factors}}
\label{sec:local-l-factors}

Let \(V\) be an \(\ell\)-adic Galois representation of
\(\operatorname{Gal}(\overline{\bQ}/\bQ)\). At a prime \(p\neq \ell\),
let \(I_p\) be inertia. The unramified part is \(V^{I_p}\). The local
\(L\)-factor is
\[
  L_p(V,s)
  =
  \det\left(1-p^{-s}\Frob_p\mid V^{I_p}\right)^{-1}.
\]
Again the determinant is the action of Frobenius on a volume line.
The Euler product
\[
  L(V,s)=\prod_pL_p(V,s)
\]
is a product of inverse determinants.

\section{Deninger-Type Determinants}
\label{sec:deninger}

For \(X/\bF_q\), Frobenius is a genuine operator on genuine
cohomology, and the zeta function is an alternating product of finite
determinants.

For arithmetic schemes over \(\Spec\bZ\), Deninger's programme asks
for an analogue:
\[
  \zeta_X(s)
  =
  \prod_i
  \det_\infty\left(s\id-\Theta\mid H^i_{\mathrm{arith}}(X)\right)^{(-1)^{i+1}}.
\]
Here \(H^i_{\mathrm{arith}}(X)\) is a conjectural arithmetic
cohomology, \(\Theta\) is a Frobenius-like infinitesimal operator, and
\(\det_\infty\) is a regularised determinant.

For \(X=\overline{\Spec\bZ}\), the expected shape is:
\[
  \xi_{\bR}(s)
  =
  \det_\infty\left(s\id-\Theta\mid
  H^1_{\mathrm{arith}}(\overline{\Spec\bZ})_{\mathrm{prim}}\right),
\]
while the cohomological degrees \(0\) and \(2\) account for the Tate
pole lines at \(s=0\) and \(s=1\).

This is not a formal analogy. It is the finite-field determinant
formula transported to the arithmetic curve. The open content is the
construction of the cohomology and operator.

\section{Arakelov Determinants}
\label{sec:arakelov-det}

Arithmetic geometry also has determinant lines with metrics. Let
\(\pi:X\to\Spec\bZ\) be an arithmetic surface and let \(E\) be a
hermitian vector bundle on \(X\). The determinant of cohomology is
\[
  \lambda(E)=\det R\pi_*E.
\]
At the archimedean fibre, analytic torsion supplies a metric on this
line, the Quillen metric. Thus the determinant line is not only
algebraic; it is Arakelov-geometric. Finite primes contribute algebraic
cohomology. The infinite place contributes spectral analysis of
Laplacians. Both are determinant constructions.

\section{The Arithmetic Lesson}
\label{sec:arithmetic-lesson}

The determinant in arithmetic geometry packages three operations:
\begin{enumerate}
\item finite-place Frobenius acting on cohomology;
\item archimedean spectral regularisation;
\item alternating cohomological signs.
\end{enumerate}
The determinant is the mechanism by which a cohomological operator
becomes a zeta or \(L\)-function.

\chapter{Geometric Representation Theory}
\label{ch:grt}

\section{The Weyl Denominator}
\label{sec:weyl-denominator}

Let \(G\) be a connected reductive group with maximal torus \(T\), Weyl
group \(W\), and positive roots \(\Phi^+\). Define
\[
  \rho:=\frac12\sum_{\alpha\in\Phi^+}\alpha.
\]
The Weyl denominator identity is
\[
  \sum_{w\in W}(-1)^{\ell(w)}e^{w\rho}
  =
  e^\rho\prod_{\alpha\in\Phi^+}(1-e^{-\alpha}).
\]
The right side is a determinant. It is the determinant of
\[
  1-\operatorname{Ad}(t)^{-1}
\]
on the positive nilpotent tangent directions:
\[
  \prod_{\alpha\in\Phi^+}(1-e^{-\alpha}).
\]
The sign \((-1)^{\ell(w)}\) is the determinant of the Weyl group
element acting on the real span of roots.

\section{\texorpdfstring{The \(GL_n\) Model}{The GLn Model}}
\label{sec:gln-model}

For \(G=GL_n\), the Weyl group is \(S_n\). Let
\[
  x_i=e^{\varepsilon_i}.
\]
The Vandermonde determinant is
\[
  \det(x_j^{n-i})_{1\leq i,j\leq n}
  =
  \prod_{i<j}(x_i-x_j).
\]
The left side is a determinant of a matrix of weights. The right side
is a product over positive roots:
\[
  \prod_{i<j}x_i(1-x_j/x_i)
  =
  e^\rho\prod_{\alpha>0}(1-e^{-\alpha})
\]
after the standard normalization. Thus the Weyl denominator is the
Vandermonde determinant in representation-theoretic clothing.

\section{Weyl Character Formula}
\label{sec:weyl-character}

For a dominant weight \(\lambda\), the Weyl character formula says
\[
  \operatorname{ch}V_\lambda
  =
  \frac{\sum_{w\in W}(-1)^{\ell(w)}e^{w(\lambda+\rho)}}
  {\sum_{w\in W}(-1)^{\ell(w)}e^{w\rho}}.
\]
This is a ratio of two alternating determinants. The denominator is
the determinant of the tangent weights of the flag variety at a fixed
point. The numerator is the same determinant twisted by the line
bundle of weight \(\lambda\).

In the Chriss-Ginzburg viewpoint, the formula is not an isolated
identity. It is the fixed-point shadow of equivariant localization on
the flag variety \(G/B\). Determinants enter because localization
divides by equivariant Euler classes, and equivariant Euler classes
are determinants of tangent representations.

\section{Localization and Euler Classes}
\label{sec:localization}

Let a torus \(T\) act on a smooth variety \(X\). At a fixed point
\(x\in X^T\), the tangent space decomposes into weight spaces:
\[
  T_xX=\bigoplus_j k_{\chi_j}.
\]
The equivariant Euler class is
\[
  e_T(T_xX)=\prod_j\chi_j.
\]
This is the determinant of the \(T\)-representation on the tangent
space, read multiplicatively in equivariant cohomology.

The Atiyah-Bott localization formula has the form
\[
  \int_X\alpha
  =
  \sum_{x\in X^T}
  \frac{\alpha|_x}{e_T(T_xX)}.
\]
The denominator is a determinant. Geometric representation theory is
full of such denominators because representation-theoretic characters
are computed by localizing sheaves on spaces with group actions.

\section{Determinant Lines for Tate Vector Spaces}
\label{sec:tate-det-lines}

Loop groups force infinite-dimensional linear algebra. Let \(V\) be a
Tate vector space, such as \(k((t))\), and let \(L,L'\subset V\) be
lattices. Their relative determinant line is
\[
  \det(L|L')
  :=
  \det\bigl(L/(L\cap L')\bigr)
  \otimes
  \det\bigl(L'/(L\cap L')\bigr)^\vee.
\]
This is finite-dimensional because the two quotient spaces are
finite-dimensional when \(L\) and \(L'\) are commensurable.

The relative determinant line measures the failure of two infinite
lattices to have the same volume. This construction produces central
extensions of loop groups. If \(g\) is a loop acting on \(V\), compare
\(L\) and \(gL\). The line
\[
  \det(L|gL)
\]
is the determinant anomaly of \(g\). Multiplication of loops gives
canonical isomorphisms only up to scalars, and those scalars form the
central extension.

\section{Affine Grassmannian and Kac-Moody Level}
\label{sec:affine-grassmannian}

For a reductive group \(G\), the affine Grassmannian is
\[
  \operatorname{Gr}_G=G((t))/G[[t]].
\]
Line bundles on \(\operatorname{Gr}_G\) encode levels of affine
Kac-Moody representations. The basic line bundle is a determinant
line. For \(G=GL_n\), a point of \(\operatorname{Gr}_{GL_n}\) is a
lattice \(L\subset k((t))^n\). Relative to the standard lattice
\[
  L_0=k[[t]]^n,
\]
the determinant line at \(L\) is
\[
  \det(L|L_0).
\]
This line is the geometric source of the basic central extension of
the loop group. The representation-theoretic level is a determinant
class.

\section{Convolution and Orientation Lines}
\label{sec:convolution}

Geometric representation theory uses correspondences:
\[
  X \xleftarrow{p} Z \xrightarrow{q} Y.
\]
A sheaf on \(X\) is pulled back to \(Z\), then pushed forward to \(Y\).
When the theory is refined to \(D\)-modules, constructible sheaves,
or coherent sheaves, the correct functor often carries shifts and
orientation lines. These orientation lines are determinant lines of
tangent or cotangent complexes:
\[
  \det(T_Z)-\det(p^*T_X)-\det(q^*T_Y)
\]
in additive notation for the Picard group.

Thus the determinant controls the normalization of convolution. Wrong
determinant lines give wrong signs, wrong shifts, and wrong dualities.

\section{Springer Theory}
\label{sec:springer}

The Springer resolution
\[
  \widetilde{\mathcal N}=T^*(G/B)\to\mathcal N
\]
is one of the central spaces of geometric representation theory. The
cotangent bundle has a canonical symplectic form. Its canonical bundle
is trivial because
\[
  \omega_{T^*X}\simeq \cO_{T^*X}
\]
for smooth \(X\). This triviality is a determinant statement: the
determinant of the cotangent directions cancels the determinant of
the tangent directions. Such cancellations are the geometric reason
Springer theory has clean duality and Weyl group actions.

\section{The Geometric Representation-Theoretic Lesson}
\label{sec:grt-lesson}

In geometric representation theory, determinants appear as:
\begin{enumerate}
\item Weyl denominators and character formula denominators;
\item equivariant Euler classes in localization;
\item determinant line bundles on affine Grassmannians;
\item central extensions of loop groups;
\item orientation lines in convolution;
\item duality normalizations in symplectic resolutions.
\end{enumerate}
Every item is the same operation: take a linearized tangent,
cotangent, lattice, or cohomology object, pass to its determinant
line, and record how symmetry acts on that line.

\chapter{The Unified Dictionary}
\label{ch:unified}

\section{One Sentence}
\label{sec:one-sentence}

The determinant is the passage from a linear operator on a space to
the induced scalar action on the space's volume line.

\section{Four Realizations}
\label{sec:four-realizations}

\begin{center}
\begin{tabular}{p{0.22\textwidth}p{0.32\textwidth}p{0.32\textwidth}}
\textbf{Field} & \textbf{Space} & \textbf{Determinant output}\\
\hline
Linear algebra & finite vector space \(V\) & scalar \(\det(T)\) on \(\det V\)\\
Algebraic geometry & vector bundle or perfect complex & determinant line bundle\\
Arithmetic geometry & cohomology with Frobenius & zeta and \(L\)-factors\\
Geometric representation theory & tangent, lattice, or sheaf-theoretic data & Euler class, denominator, central extension, level\\
\end{tabular}
\end{center}

\section{The Three Tests}
\label{sec:three-tests}

When a claimed determinant appears, test it by three questions.

\begin{enumerate}
\item What is the underlying linear or derived object?
\item What is its determinant line?
\item Which operator, symmetry, or Frobenius acts on that line?
\end{enumerate}

If these three objects are not named, the determinant has not yet been
constructed. It has only been notated.

\section{The Arithmetic Warning}
\label{sec:arithmetic-warning}

For a finite-dimensional operator, the determinant is automatic. For
an infinite-dimensional arithmetic operator, it is not. A formula
\[
  \xi(s)=\det_\infty(s\id-\Theta)
\]
requires:
\begin{enumerate}
\item a cohomological space;
\item a densely-defined operator \(\Theta\);
\item a spectral theory for \(\Theta\);
\item a regularised determinant;
\item a trace formula linking the determinant to primes and
archimedean terms.
\end{enumerate}

The finite-field formula proves that this is the right dream. It does
not construct the number-field object. The construction of that object
is the substance of a Deninger-type arithmetic cohomology.

\chapter*{Closing}
\addcontentsline{toc}{chapter}{Closing}

The determinant begins as the answer to a baby question:

\[
  \text{by what scalar does }T\text{ multiply volume?}
\]

The rigorous answer is:
\[
  \bigwedge^nT=\det(T)\cdot\id_{\bigwedge^nV}.
\]

Algebraic geometry globalizes the volume line into a determinant line
bundle. Arithmetic geometry makes Frobenius act on cohomological
volume lines and obtains zeta functions. Geometric representation
theory makes symmetry act on tangent, lattice, and sheaf-theoretic
determinant lines and obtains denominators, levels, Euler classes, and
central extensions.

The same construction survives every change of field:
\[
  \boxed{\text{determinant}=\text{action on the top exterior line}.}
\]

\end{document}
